\documentclass[../main.tex]{subfiles}
\ifSubfilesClassLoaded{\setcounter{chapter}{10}}{}
\begin{document}

\chapter{Klastering: K-Means dan Silhouette}

\begin{subcpmk}
  \item \textbf{Sub-CPMK-P8:} K-Means pada fitur numerik ter-skala; argumen pemilihan $k$ dengan silhouette atau inertia (RPS Mg.\ 10).
\end{subcpmk}

\SesiRPS{10}{P8}{$\sim$170 menit}{CPMK-P3}

\section{Teori}

\begin{konsep}
Klastering \textit{unsupervised}: tidak ada label. K-Means meminimalkan inertia dalam cluster; hasil dipengaruhi skala fitur---gunakan \texttt{StandardScaler}.
\end{konsep}

\section{Contoh kode}

\begin{lstlisting}[caption=Data dan scaling]
from sklearn.datasets import make_blobs
from sklearn.preprocessing import StandardScaler

X, _ = make_blobs(n_samples=400, centers=4, cluster_std=0.6, random_state=42)
scaler = StandardScaler()
Xs = scaler.fit_transform(X)
\end{lstlisting}

\begin{lstlisting}[caption=KMeans dan inertia]
from sklearn.cluster import KMeans
km = KMeans(n_clusters=4, random_state=42, n_init="auto")
labels = km.fit_predict(Xs)
print("inertia", km.inertia_)
\end{lstlisting}

\begin{lstlisting}[caption=Silhouette score untuk beberapa k]
from sklearn.metrics import silhouette_score
for k in range(2, 8):
    kmk = KMeans(n_clusters=k, random_state=42, n_init="auto")
    lab = kmk.fit_predict(Xs)
    s = silhouette_score(Xs, lab)
    print(k, s)
\end{lstlisting}

\begin{lstlisting}[caption=Profil cluster (rata-rata per fitur)]
import numpy as np
import pandas as pd
dfp = pd.DataFrame(Xs)
dfp["cluster"] = labels
print(dfp.groupby("cluster").mean())
\end{lstlisting}

\section{Referensi}

\begin{itemize}[leftmargin=*]
  \item sklearn unsupervised. \url{https://scikit-learn.org/stable/unsupervised_learning.html}
  \item Contoh silhouette K-Means. \url{https://scikit-learn.org/stable/auto_examples/cluster/plot_kmeans_silhouette_analysis.html}
\end{itemize}

\begin{asesmen}
\textbf{Tugas modul:} pemilihan $k$ berargumen + profil deskriptif cluster.
\end{asesmen}

\begin{rangkuman}
Minggu 10 mengeksplorasi pola tanpa label dengan K-Means.

\textbf{Kata kunci:} KMeans, silhouette, inertia
\end{rangkuman}

\end{document}
